\documentclass[11pt]{article}

\usepackage{maiacustom}

\begin{document}

\psettitle{Astronomy Question Bank}

These questions were carefully written or selected to prepare students for the astronomy team selection process in Brazil. Some questions are not original; those are marked with () before the statement. The question-bank template is the same one used by Professor Kevin Zhou. His work is invaluable, and many ideas in this set can be found in his handouts.

% \begin{psidea}{Idea title}{}
% Idea
% \end{psidea}

% \begin{psexample}{Example title}{}
% Example
% \end{psexample}

% \begin{pssolution*}{}{}
% Solution
% \end{pssolution*}

% \begin{psremark*}{Remark title}{}
% Remark
% \end{psremark*}
\section{Positional Astronomy}

\pts{5} % problem of the week 114
\begin{pproblem} After getting lost on a voyage, you and Professor Klafke are shipwrecked on an island and find the following sundial:
    \begin{figure}[H]
        \centering
        \includegraphics[width=0.5\linewidth]{imagens/q6.jpg}
        \caption{Sundial}
    \end{figure}
    It is now your task to discover properties of the sundial.
    \begin{enumerate}[label=\textbf{\alph*)}]
        \item In which hemisphere did you find yourselves? How did you reach that conclusion?
        \item After several weeks on the island, you notice that the sundial's shadow seems to follow a straight line. Letting \(\phi\) be the latitude at which you found yourselves, determine the Sun's declination on that day and explain why this happens on only \(2\) days of the year.
        \item Prove that, when the Sun has the declination found above, the shadow follows a straight line. This can be done by the following steps:
        \begin{enumerate}[label=\roman*)]
            \item Determine the orientation of the line and explain why it must have that orientation;
            \item Determine the length of the shadow for a given position of the Sun in alt-azimuth coordinates;
            \item Derive a relation between altitude and azimuth, given that the tip of the shadow lies on the line;
            \item Determine a constant quantity and show that it is constant for every position of the Sun on that day.
        \end{enumerate}
    \end{enumerate}
\end{pproblem}

    \pts{5}
\begin{pproblem} (Adapted from T1 - 2024)
    The equation of time is the difference between the right ascension of the mean Sun and the right ascension of the true Sun. With that in mind, let us calculate some of its properties.
    \begin{alternativas}
        \item Find an expression for the equation of time, neglecting Earth's eccentricity, i.e.: the Sun moving at constant speed along the ecliptic. Leave your answer in terms of the time since the March equinox \(T_M\), Earth's period \(T\), and the obliquity of the ecliptic \(\epsilon\).
        \item Now, taking Earth's eccentricity into account, the equation of time has the form:
        \[E.T. = -2e \cdot \sin(k_1 \cdot M) + \tan^2 \left( \frac{\epsilon}{2} \right) \cdot \sin\left( k_2 \cdot (M + \lambda_P) \right)\]
        where \(M\) is the mean anomaly, \(\lambda_P\) the ecliptic longitude of periapsis, and \(e\) Earth's eccentricity. Determine the constants \(k_1\) and \(k_2\). Think about the symmetries involving each term of the equation.
        \\
        The aim of the question is now to find the Sun's declination at the node of the analemma. See the following figure:
        \begin{figure}[H]
            \centering
            \includegraphics[width=0.7\textwidth]{imagens/q16.png}
            \caption{Analemma diagram}
        \end{figure}
        The node is the point at which the path crosses itself.
        \item For small values of the Sun's declination, obtain a formula for it in terms of \(M\), \(\lambda_P\), and \(\epsilon\). Also assume that the eccentricity and the orbital inclination are sufficiently small.
        \item Find the two mean anomalies that have the same declination and the same \(E.T.\).
        \item Find the value of the solar declination at the node.
    \end{alternativas}
    \end{pproblem}

    
    \pts{5}
    \begin{pproblem}
        Porílio, after a long day of classes at ITA, on the summer solstice (\(\phi = 23^\circ 11' S, \ \lambda = 45^\circ 53' W\)) wants to watch the sunset. On the day in question, the equation of time is \(E.T. = -3\) minutes. Porílio does not want to miss the Sun under any circumstances and starts running so that the (center of the) Sun stays on the horizon.
    
        \begin{alternativas}
            \item São José dos Campos follows Brasília time (UTC-3h). At what civil time should Porílio start his run?
            \item If Porílio runs toward one of the geographic poles, what should his initial speed be?
            \item If Porílio decides to climb a mountain of inclination \(i = 30^\circ\) at a constant speed \(v = 1\) m/s, what is the longest time for which Porílio can keep the Sun above the horizon?
            \textbf{Hint: } Small-angle approximations may be useful.
        \end{alternativas}
\end{pproblem}

\pts{3}
\begin{pproblem}
    Mauí, a renowned astronomer, wants to spend his year-end vacation in Sergipe \((10^\circ \ 54'\ 33'' S, \ 37^\circ 4'29'' W)\) and wants to know the rising time of one of his favorite stars, All Kaf al Dij Ma \Romannum{3} \((\delta = +10^\circ \ 56'\ 58'', \ \alpha = 2h \ 58m \ 43s\)), and needs your help to calculate its rising time in the following situations:
    
    \begin{alternativas}
        \item Calculate the rising time of All Kaf al Dij Ma \Romannum{3} at the December solstice.
        \item How long will the star spend above the horizon?
        \item Estimate the length of time during which the star is visible.
    \end{alternativas}

    \end{pproblem}

    \pts{4}
\begin{pproblem}(Lucas Cavalcante - \href{https://noic.com.br/olimpiadas/astronomia/problemas-da-semana/astronomia-semana-111/}{Week 111})
    Consider a horizontal coordinate system in which azimuth \(0\) corresponds to the north cardinal point. One asteroid was observed at altitudes \(h_1\) and \(h_2\) and azimuths \(A_1\) and \(A_2\), while another asteroid was observed at the coordinates \(h_3\), \(h_4\), \(A_3\), and \(A_4\). Find the coordinates of the points where the asteroids meet.
    
    \begin{psidea}{Einstein notation}{}
    While solving this problem, I will use Einstein notation. It is a tool from physics that makes vector calculations easier. In ordinary notation, a vector is written as
    \[\mathbf{a} = (a_1 \hat{\mathbf{e}}_1, \ a_2 \hat{\mathbf{e}}_2, \ a_3 \hat{\mathbf{e}}_3)\]
    In Einstein notation, this is replaced by:
    \[\mathbf{a} = a_i \hat{\mathbf{e}}_i\]
    Repeated indices indicate a sum (\(a_i \hat{\mathbf{e}}_i = a_1 \hat{\mathbf{e}}_1 + a_2 \hat{\mathbf{e}}_2 + a_3 \hat{\mathbf{e}}_3 + ...\))
    
    The scalar product of two vectors, in Einstein notation, is defined by:
    \[\mathbf{a}\cdot \mathbf{b} = a_i b_j \delta_{ij}\]
    where \(\delta_{ij}\) is the \textit{Kronecker delta}, defined by:
    \[\delta_{ij} = \left \{ \begin{matrix} 0, & \mbox{if } i \ne j \\ 
                                            1, & \mbox{if } i = j\end{matrix} \right.\]
    
    Thus, \(\mathbf{a}\cdot \mathbf{b} = a_i b_i \equiv a_1 b_1 + a_2 b_2 + a_3 b_3\)

    For the cross product, we have:
        \[
        \mathbf{a} \times \mathbf{b} = a_i b_j \Epsilon_{ijk} \hat{\mathbf{e}}_k
        \]
        Here, \(\Epsilon_{ijk}\) is called the \textit{Levi-Civita tensor}, or Levi-Civita symbol, and is defined as:
    
        \[
        \Epsilon_{ijk} = 
        \begin{cases} 
        +1 & \text{if } (i, j, k) \text{ is an even permutation of } (1, 2, 3), \\
        -1 & \text{if } (i, j, k) \text{ is an odd permutation of } (1, 2, 3), \\
        0 & \text{if two or more indices are equal}.
        \end{cases}
        \]

        For example, \(\Epsilon_{123} = \Epsilon_{312} = \Epsilon_{231} = 1\) and \(\Epsilon_{132} = \Epsilon_{213} = \Epsilon_{321} = -1\).
        Explicitly 
        \[\mathbf{a} \times \mathbf{b} = (a_2 b_3 - a_3 b_2) \hat{\mathbf{e}}_1 + (a_3 b_1 - a_1 b_3) \hat{\mathbf{e}}_2 + (a_1 b_2 - a_2 b_1) \hat{\mathbf{e}}_3\]
    \end{psidea}
\end{pproblem}


\pts{5}
\begin{pproblem} (Problem set \(1\) - Vinhedo 2024) Consider a sundial consisting of a vertical dial
    and a gnomon, as shown in the figure.
    \begin{figure}[H]
        \centering
        \includegraphics{imagens/q20.png}
        \caption{Sundial diagram.}
    \end{figure}
    \begin{alternativas}
        \item Suppose a vertical sundial whose gnomon is correctly pointed south.
        Let \(\phi\) be the observer's latitude, \(H\) the Sun's hour angle, and \(\delta\) the solar declination. Determine
        the relation between \(\theta\) and \(\gamma\) in terms of these variables.
        
        \item From the previous item, what must \(\theta\) be for the sundial to work throughout the year?

        \item Using the correct value of \(\theta\), determine the relation between \(\gamma\) and \(H\), in terms of \(\phi\) only.
    \end{alternativas}
\end{pproblem}


\pts{2} 
\begin{pproblem}
    What is the smallest height a giant (who lives a VERY long time), located at the South Pole, must have in order to see every star in the sky?
\end{pproblem}


\pts{3}
\begin{pproblem}
    Toleduardo wants to watch the sunset, but he spent too long at SorveterITA. Toleduardo is a \textit{Just in time} man and wants to know how long the sunset will last, so that he can be late at his leisure. Given that Toleduardo only goes to SorveterITA on March 21 or September 22, calculate the duration of the sunset at SorveterITA.
    \\
    \textbf{Data: } Latitude of SorveterITA \(\phi = 23^\circ \ 10' \ 45''S, \ \ \lambda = 45^\circ 53' 14''W\).
\end{pproblem}


\pts{3}
\begin{pproblem} (Problem set 2 - Vinhedo 2021)
    In February 2015, the Moon began a cycle of monthly occultations of Aldebaran ($\alpha$ Tau). That is, every month the Moon passed in front of Aldebaran for an observer on Earth. Note that these occultations did not necessarily occur for observers at the same location every month. 

    Calculate the date (month and year) when this cycle of monthly occultations ends. Assume that the Moon's orbit is circular.

    \subsection*{Data:}

    \begin{itemize}
        \item Ecliptic latitude of Aldebaran = $5.47^\circ$
        \item Nodal precession period of the Moon = 18.6 years
    \end{itemize}
\end{pproblem}


\pts{4}
\begin{pproblem} (Problem set 2 - Vinhedo 2021)
    Miguel lives on an isolated island in the South Pacific Ocean, at longitude
    $\lambda = 176^{\circ}09^{\prime}137.7^{\prime\prime}\text{W}$. Over the year, the place where the Sun rises, as seen by Miguel, varies by $\Delta A = 67^{\circ}03^{\prime}81^{\prime\prime}$ along the horizon. For the first two items, neglect atmospheric refraction. With this information, find:

    \begin{alternativas}
        \item The latitude $\phi$ and the name of the island. (Consult Google Earth or similar software)
        \item The range of times at which the Sun rises on the island, given the peculiar time zone $UT + 12\frac{3}{4}$.
        \item Taking atmospheric refraction into account, would the length of the interval in the previous item decrease, increase, or stay constant?
    \end{alternativas}

\end{pproblem}


\pts{3}
\begin{pproblem}(Problem set 2 - Vinhedo 2022) 
    Bruno decided to rent a house for his vacation in Cuiabá ($15.3^{\circ} \text{S}$, $56.1^{\circ} \text{W}$). As a good astronomer, Bruno spent his nights sitting in a chair, watching the stars through a giant door facing the south cardinal point.

    The door was 4.00 meters high and 1.50 meters wide. Bruno used to sit
    1.00 meter from the door, perfectly aligned with its center horizontally. That is, the line
    segment between Bruno's eyes and the point exactly halfway across the door horizontally and
    at eye level forms an angle of $90^{\circ}$ with the plane of the door. Bruno's eyes are
    1.20 meters above the floor when he is sitting in the chair.
    
    To make his observations easier, Bruno created a coordinate system based on the position of the
    door, from which he saw the stars at the place where he was sitting, using meters as
    the reference unit. The origin of the system is at the lower-left corner. The coordinates
    in $x$ increase to the right and the coordinates in $y$ increase upward. Thus, a
    star seen 1 meter in from the left side of the door and 2 meters above the floor would be represented
    by the coordinates $(1, 2)$.
    
    Bruno was very interested in Shaula ($\lambda \ \text{Sco}, \delta = 37.1^{\circ} \text{S}$). Determine the coordinates of
    Shaula at the instant the star became visible to Bruno when observed through the door. Assume that Shaula was below the horizon when Bruno began watching the sky.
\end{pproblem}


\end{document}
